\documentclass[10pt]{article}

% Specify the margins and text size.
\setlength{\textwidth}{6.4in}
\setlength{\textheight}{9.5in}
\setlength{\oddsidemargin}{0pt}
\setlength{\evensidemargin}{0pt}
\setlength{\topmargin}{0pt}
\setlength{\hoffset}{.05in}
\setlength{\voffset}{-.8in}

\setlength{\parskip}{5pt}
\setlength{\parindent}{0pt}

% Load some fonts and symbol packages
\usepackage{latexsym}
\usepackage{pifont}       % contains 'star' symbol for counterinsurgency handbook title
\usepackage{yfonts} 
\usepackage{amsmath}
\usepackage{amsfonts}

\usepackage{graphicx}     % actually, this is loaded by pstricks
\usepackage[T1]{fontenc}
\usepackage{ifthen}
\usepackage{pstricks,pst-grad,pst-text,pst-node,multido,pst-plot,calc,pst-3dplot}
%\usepackage[all]{xy}
%\usepackage{animate}

% The hyperref package inserts links into the pdf.
\definecolor{MyLinkColor}{rgb}{.1,.2,1}
\definecolor{MyCiteColor}{rgb}{.1,1,.2}
\definecolor{MyURLColor}{rgb}{.4,.4,.4}
\usepackage[backref=true,pagebackref=false,hyperindex,colorlinks=true,
  linkcolor=MyLinkColor,urlcolor=MyURLColor]{hyperref}


% The tweaklist package is something I found on the web.  It provides a simple interface
% for making changes to spacing used in the itemize and enumerate environments.  Comment
% this out if you don't care to use tweaklist.
\usepackage{tweaklist}
\renewcommand{\itemhook}{\setlength{\parskip}{2pt}\setlength{\parsep}%
{1pt}\setlength{\topsep}{0pt}\setlength{\itemsep}{0pt}}

\newcommand{\U}{\underline{\hspace{5pt}}}

\usepackage{listings}
\newcounter{EX}\setcounter{EX}{1}
\newcommand{\EXERCISE}{\arabic{EX}.\stepcounter{EX} }

\begin{document}
\pagestyle{empty}
\lstset{language=R, showspaces=false, showstringspaces=false}

\href{http://www.su.edu}{\includegraphics[height=1.75cm]{sulogo.eps}}
\vspace{-1.69cm}

{\small \hfill
\begin{tabular}{cl}
& Math 207\\& Introduction to Statistics\\
\end{tabular}
}
\setlength{\baselineskip}{1.05\baselineskip}
\medskip
\medskip

\begin{center}
\textbf{\large  Variables}
\end{center}

\newcommand{\SUBX}{\smallskip\hspace{10pt}}
\newcommand{\BSK}{\vspace{.14in}}

A (random) \textbf{variable} is a measurement that depends on the outcome 
   of a (random) event.
\begin{itemize}
\item \textbf{Quantitative} variables have numeric values.  
  \begin{itemize} 
     \item \textbf{Continuous} variables can assume a continuum of values:  Examples include 
      income, temperature, pressure, mass, and speed. 
     \item A \textbf{discrete} variable can assume only finitely (or countably) many values.
       Examples include: family size, and number of engine cylinders.\vspace{2pt}
  \end{itemize}
\item \textbf{Qualitative} variables are non-numeric.  
   \begin{itemize}
      \item A qualitative variable can be \textbf{ordered} such as customer 
          satisfaction (good, better or best)
      \item or \textbf{unordered} such as eye color or automobile transmission type.
   \end{itemize}
\end{itemize}

When working with a variable, it is important to understand what kind of variable
it is.
The kind of variable we are working with influences the statistical things that we 
can do with it.
We can compute averages, standard deviations and other numeric statistics for 
quantitative variables.  We can't do this for qualitative variables.
We can make histograms of any kind of variable.  If the variable is qualitative 
but unordered, then the histogram boxes can be placed in any order.
Variable type also determines the kinds of computer data structures used to store its
values.\vspace{-5pt}
\begin{center}
{\footnotesize\begin{pspicture}(-1.5,0)(7.5,2.2)
\psset{linewidth=0.02, nodesep=2pt, xunit=3}
\rput(1,2){\rnode{Variables}{Variables}}
\rput(0,1){\rnode{Qualitative}{Qualitative}}
\rput(2,1){\rnode{Quantitative}{Quantitative}}
\rput(1.5,0){\rnode{Discrete}{Discrete}}
\rput(2.5,0){\rnode{Continuous}{Continuous}}
\rput(-0.5,0){\rnode{Ordered}{Ordered}}
\rput(0.5,0){\rnode{Unordered}{Unordered}}
\ncline{Variables}{Qualitative}
\ncline{Variables}{Quantitative}
\ncline{Quantitative}{Discrete}
\ncline{Quantitative}{Continuous}
\ncline{Qualitative}{Ordered}
\ncline{Qualitative}{Unordered}
\end{pspicture}}
\end{center}


1. Classify each of the following variables as qualitative or quantitative;  
if quantitative, as discrete or continuous;  if qualitative, as ordered or
unordered.\vspace{-5pt}
\begin{center}
\begin{tabular}{lll}
(a) occupation & (b) region of residence & (c) weight\\[5pt]
(d) height     & (e) number of automobiles owned\span
\end{tabular}
\end{center}

2. The following variables 
are in  the \texttt{adult.csv} file of the NHANESIII data set.
Classify each as qualitative or quantitative;  
if quantitative, as discrete or continuous;  if qualitative, as ordered or
unordered.\vspace{-5pt}
\begin{center}
{\setlength{\tabcolsep}{20pt}\begin{tabular}{llll}
(a) \texttt{DMARACER} & (b) \texttt{HSAGEIR} & (c) \texttt{MXPLANG} & (d) \texttt{HFA8R}\\[5pt]
(e) \texttt{HFA13}    & (f) \texttt{HFB1}    & (g) \texttt{HAB6S}   & (h) \texttt{SEQN}\\
\end{tabular}}
\end{center}

3. The following variables 
are in  the \texttt{exam.csv} file of the NHANESIII data set.
Classify each as qualitative or quantitative;  
if quantitative, as discrete or continuous;  if qualitative, as ordered or
unordered.\vspace{-5pt}
\begin{center}
{\setlength{\tabcolsep}{20pt}\begin{tabular}{llll}
(a) \texttt{BMPWT} & (b) \texttt{BMPHT} & (c) \texttt{PEP5R} & (d) \texttt{PEPPREG}\\[5pt]
(e) \texttt{DMPCREGN}    & (f) \texttt{DMPPIR}  & (g) \texttt{DEPCS2}  & (h) \texttt{SEQN}\\
\end{tabular}}
\end{center}

4. The following variables 
are in  the \texttt{lab.csv} file of the NHANESIII data set.
Classify each as qualitative or quantitative;  
if quantitative, as discrete or continuous;  if qualitative, as ordered or
unordered.\vspace{-5pt}
\begin{center}
{\setlength{\tabcolsep}{20pt}\begin{tabular}{llll}
(a) \texttt{PHPCHM2} & (b) \texttt{PHPSNDA} & (c) \texttt{WCP} & (d) \texttt{RCP}\\[5pt]
(e) \texttt{PBP}    & (f) \texttt{LCP}  & (g) \texttt{UIP}  & (h) \texttt{SEQN}\\
\end{tabular}}
\end{center}




\vfill
\eject

\begin{center}
\textbf{\large Solutions}
\end{center}

1. (a) occupation is a qualitative variable.  It's not a number.
(b) region of residence is also qualitative.  It's not a number either.
(c) weight is quantitative.  It's a continuous value (can be measured to many decimal places 
if you have the right equipment).
(d) height is quantitative and continuous just like weight is.
(e) number of automobiles owned is quantitative and discrete (assume you can't 
own fractions of a car!).
\bigskip

2. 

(a) \texttt{DMARACER} is an unordered, qualitative variable that gives information
about the subject's race.

(b) \texttt{HSAGEIR} is a discrete, quantitative variable that gives the age
of the subject in years.

(c) \texttt{MXPLANG} is an unordered, qualitative variable that gives information
about the language used in the interview.

(d) \texttt{HFA8R} is a discrete, quantitative variable that gives information
about the number of years of school completed by the subject.  Blank and Don't know
are two possible responses that need to be omitted to compute statistics.

(e) \texttt{HFA13} is an unordered, qualitative variable that gives information
about whether or not the subject ever served in the U.S. armed forces.

(f) \texttt{HFB1}  is an unordered, qualitative variable that gives information
about whether or not the subject is covered by Medicare.

(g) \texttt{HAB6S}  is a discrete, quantitative variable that gives the number
of months since the subject last saw or talked to a medical doctor or other 
health professional.

(h) \texttt{SEQN} is the identification number of the subject.  Typically this would
be viewed as an ordered, qualitative variable.
\medskip


3. (a) \texttt{BMPWT} is weight in kg.  If we eliminate the blank responses, this is a 
continuous, quantitative variable.  

(b) \texttt{BMPHT} is standing height in cm.  
It is a  continuous, quantitative variable. 

(c) \texttt{PEP5R} is central pulse rate
in beats per minute.  It is a discrete, quantitative variable.

(d) \texttt{PEPPREG} is an unordered, qualitative variable indicating 
whether or not the subject is pregnant.  Blank and Don't know are also possible responses.

(e) \texttt{DMPCREGN}  is the census region (NW, MW, S, W) of the subject.  It is 
   an unordered, qualitative variable.

(f) \texttt{DMPPIR}  is the subject's poverty income ratio.  It is a continuous,
  quantitative variable.

(g) \texttt{DEPCS2}  gives unordered, qualitative  information about a the upper left
central incisor tooth.  

(h) \texttt{SEQN} is the identification number of the subject.  Typically this would
be viewed as an ordered, qualitative variable.
\medskip

4.

(a) \texttt{PHPCHM2} is an unordered, qualitative variable indicating whether
or not the subject has received chemotherapy treatment within the past four weeks.

(b) \texttt{PHPSNDA} is an ordered, qualitative variable indicating the day that
the participant last ate.

(c) \texttt{WCP} is white blood cell count.
It is a continuous,  quantitative variable.

(d) \texttt{RCP} is red blood cell count.  
It is a continuous,  quantitative variable.

(e) \texttt{PBP} is lead level in the blood sample measured in ug/dL.  
It is a continuous,  quantitative variable.

(f) \texttt{LCP} is serum LDL cholesterol measured in mg/dL.
It is a continuous,  quantitative variable.

(g) \texttt{UIP}  is urinary iodine measured in ug/dL.  
It is a continuous,  quantitative variable.

(h) \texttt{SEQN} is the identification number of the subject.  Typically this would
be viewed as an ordered, qualitative variable.
\vfill
\eject
\end{document}
